\cotefile{162-5}
\scanpage{34}
On a aussi un foncteur évident
\[ T: (\mathrm{Ssimpl}) \longrightarrow (\mathrm{Cat}) , \]
en associant à tout ensemble semi-simplicial X la catégorie $T(X) = \Delta_{/X}$
des simplexes sur X, dont l'ensemble des objets est la réunion d\deleted{e}isjointe
des $X_n$ .... (c'est une catégorie fibrée sur la catégorie $\unsure{\Delta|\underline{\Delta}}$ des simplexes
types, à fibres les catégories discrètes définies par les $X_n$). Ceci posé,
Quillen prouve que pour tout X, ST(X) est isomorphe canoniquement à X dans
la catégorie homotopique construite avec (Ssimpl), et que pour toute C,
la catégorie TS(C) est \inserted{canoniquement} ``homotopiq\deleted{y}uement équivalente à C'' \deleted{dans le sens} i.e.
\inserted{\deleted{\unsure{à|è}}ca\deleted{\unsure{C}}noniquement} isomorphe à C dans la catégorie quotient de (Cat) obtenue en inversant les
foncteurs qui sont des homotopismes. \deleted{Il en résulte} Ces isomorphismes sont
fonctoriels en X. I\={l} en résulte formellement qu'un morphisme $f:X\to Y$ dans
(Ssimpl) est un homotopisme si et seulement si il en est ainsi de $T(f):T(X)$
$\to T(Y)$, d'où des foncteurs $S'$: $(\mathrm{Cat})' \longrightarrow (\mathrm{Ssimpl})'$ et $T'$:$(\mathrm{Ssimpl})'\longrightarrow(\mathrm{Cat})'$
entre les catégories ``homotopiques'', \inserted{construites avec (Cat) resp,(Ssimpl),} qui sont quasi-inverses l'un de l'aut-\\
re.

\deleted{\unsure{d}}De plus, Quillen construit un isomorphisme canonique et fonctoriel
dans (Cat)' entre C et la catégorie opposée $C^{\mathrm{o}}$ , ou ce qui revient au
m\deleted{\unsure{m}}ê\deleted{\unsure{e}}me, un isomorphisme canonique et fonctoriel dans (Ssimpl)' entre S(C) et
S($C^{\mathrm{o}}$). \deleted{Cela lui permet alors, pour} La définition est telle que le
foncteur induit sur les systêmes locaux sur C tra\deleted{\unsure{b|h}}ns\deleted{o}f\deleted{f}orme le foncteur
contrav\deleted{\unsure{r}}ariant F sur C, transformant t\deleted{\unsure{s|e}}oute flêche en flêche inversible, en
le foncteur covariant (i.e. contravariant sur $C^{\mathrm{o}}$) ayant mêmes valeurs sur
les objets, et obtenu sur les flêches en remplaçant F(u) par $F(u)^{-1}$; en
d'autres termes, l'effet de l'homotopisme de Quillen sur les group\deleted{i}oides
fondamentaux est l'isomorphisme évident entre les groupoides fondamentaux
de C et de $C^{\mathrm{o}}$ ,\deleted{qui est \unsure{l'ideni|l'idni}} compte tenu que le deuxième est l'opposé du

\margin{33}
